\scanpage{20}
\uline{Cohomologie des quadriques}

$Q^n$ quadrique non singulière de dim $n$ sur $k$ alg. clos.

$Q^{n-1}$ section à l'infini.
\[
U^n=AQ^n=Q^n-Q^{n-1}
\]
\margin{pour $A$ coefficient annulé par $m$}
\[
H_c^i(U^n,A)\simeq
\left\{\begin{array}{ll}
0&\text{si }i\neq n,2n\\
A\otimes\unsure{V'_n}&\text{si }i=n\\
A\otimes\mu_m(-n)&\text{si }i=2n
\end{array}\right.
\quad\inserted{\text{si }n\neq0}
\]
\margin{\[
\unsure{V'_n}=\check V_n\otimes\mu_m(-n)=
\begin{cases}
\mu_m(-\frac n2)&n\text{ pair}\\
\mu_m(\unsure{-\frac{n+1}{2}})&n\text{ imp.}
\end{cases}
\]
\[
V'_n=V_n=\mu_m(-\tfrac n2)\quad n\text{ pair}
\]
\[
V'_n=V_n\otimes\mu_m(1)\quad n\text{ impair}
\]}
dualement
\diagram{\begin{tikzpicture}
\node at (-2,0) {$H^i(U^n,A)\simeq$};
\node at (0,0.65) {$0$}; \node[anchor=west] at (1,0.65) {si $i\neq0,n$};
\node (a) at (0,0) {$A\otimes V_n$}; \node[anchor=west] at (1,0) {$i=0$};
\node (b) at (0,-0.7) {$A$}; \node[anchor=west] at (1,-0.7) {si $i=n$.};
\draw (-0.7,1) .. controls (-0.9,0.9) and (-0.9,0.2) .. (-1,0.1) .. controls (-0.9,0) and (-0.9,-0.8) .. (-0.7,-1);
\draw[->] (a.west) to[bend right=30] (b.west);
\draw[->] (b.west) to[bend left=50] (a.west);
\node[draw,ellipse] at (-0.8,-0.85) {si $n\neq0$};
\end{tikzpicture}}
\margin{\[
V_n\simeq\mu_m(-\tfrac n2)\ \text{si }n\text{ pair},\qquad
V_n\simeq\mu_m(\unsure{\tfrac{-n+1}{2}})\ \text{si }n\text{ impair}
\]}
\margin{\unsure{Je m'y perds, ces deux sont définis à un signe près}\ldots
\figure{\begin{tikzpicture}\draw[<->,very thick,double] (0,0)--(0,1.5);\end{tikzpicture}}}

$H_c^i(U^n,A)\overset{\varphi}{\longrightarrow}H_{\unsure{c}}^i(\unsure{U^n|W^n},A)$ est \uline{nul} \uuline{sauf} si $n$ pair, $i=n$, dans \unsure{ce cas} devient un isomorphisme \margin{\uline{exc. 2}}. Il y a un gén. $\ell$ de $H_c^i(U^n,A)$ tel que $\langle\ell,\varphi(\ell)\rangle=2$.
\[
\deleted{\unsure{H^i(Q^n)\subset H^i(Q^n)}}
\]
\[
H^{i-1}(U^n)\text{---}H^{i-2}(Q^{n-1})(-1)\longrightarrow H^i(Q^n)
\longrightarrow H^i(\unsure{U^n})\text{---}H^{i-1}(Q^{n-1})(-1)\text{---}\cdots
\]
\[
H^{i-2}(Q^{n-1})(-1)\overset{\alpha_*}{\longrightarrow}H^i(Q^n)
\text{ est }\left\{\begin{array}{ll}
\text{surjectif}&\text{si }i\neq0,n\\
\text{injectif}&\text{si }i\neq1,n+1
\end{array}\right.
\quad
\left|\begin{array}{l}
\text{bijectif si}\\
i\neq0,\unsure{1},n,n+1
\end{array}\right.
\]
Que se passe-t-il pour $i=0,1,n,n+1$ ??
    
Pour $i=0,1$ $\alpha_*$ est \uline{nul}
\margin{[$i=1$, \deleted{\unsure{bij}} $\alpha$ bijectif, car les deux \unsure{mbs} nuls)}
    
Pour $i=n,n+1$ on a la suite exacte \margin{(si $\boxed{n\neq1}$ donc $n-1\neq0$]}
\[
0\text{---}H^{n-2}(Q^{n-1})(-1)\longrightarrow H^n(Q^n)
\longrightarrow H^n(U^n)\text{---}H^{n-1}(Q^{n-1})(-1)\text{---}H^{n+1}(Q^n)\text{---}0
\]
1) $n$ \uline{pair} $>0$, la suite exacte se réduit :
\[
0\longrightarrow H^{n-2}(Q^{n-1})(-1)\longrightarrow H^n(Q^n)\longrightarrow H^n(U^n)\text{---}0
\]
donc $H^n(Q^n)$ est extension de $H^n(U^n)$ par $H^{n-2}(Q^{n-1})(-1)$ \ldots
\mspage{19}
